\clearpage\thispagestyle{empty}
\pdfbookmark[0]{Quellen zu Lindenaus Bericht über den Sonnendurchmesser von 1810}{lindenausolar1810-sources}
\begin{center}\textbf{Quellen zu Lindenaus Bericht über den Sonnendurchmesser von 1810}\end{center}
\noindent Diese KI-übersetzte Fassung in modernem Deutsch folgt Carl Friedrich Gauß, \href{https://gdz.sub.uni-goettingen.de/id/PPN236007467}{\emph{Werke}, Band~VI} (Göttingen, 1874), S.~525–526. Der Bericht erschien zuerst ohne Verfassernamen unter Göttingen in den \emph{Göttingischen gelehrten Anzeigen}, Stück~95, 16.~Juni 1810, \href{https://images.p.sub.uni-goettingen.de/iiif/image/gdz:PPN319721507_1810:00000969/full/1650,/0/default.jpg}{S.~937}–\href{https://images.p.sub.uni-goettingen.de/iiif/image/gdz:PPN319721507_1810:00000972/full/1650,/0/default.jpg}{940}. Bernhard von Lindenau teilte Gauß seine neuen Ergebnisse in einem Schreiben vom 23.~Mai mit. Der Bericht verweist auf den \hyperlink{native-524-lower}{Bericht vom 17.~Juli 1809} und auf die \hyperlink{native-518}{in der Werkausgabe angeführte Formel}. Auch der \href{https://images.p.sub.uni-goettingen.de/iiif/image/gdz:PPN319721507_1810:00000970/full/1650,/0/default.jpg}{Erstdruck des Berichts von 1810} nennt 0,9631 als Druckfehler für 0,09631. Übersetzung und Satz mit KI erstellt.
\clearpage
